% Bachelorarbeit ---------------------------------------------------------------
%   Diese Datei gliedert die Bachelorarbeit und setzt Grundlagen
% ------------------------------------------------------------------------------
\documentclass[
    11pt,                               % Schriftgröße
    a4paper,                            % Papierformat DIN
    titlepage,                          % es wird eine Titelseite verwendet
    parskip=half,                       % Abstand zwischen Absätzen (halbe Zeile) oder full (ganze Zeile)
    headings=normal,                    % Größe der Überschriften verkleinern
    leqno,                              % Formeln werden mittig dargestellt und erhalten Nummer
    twoside,                            % Zweiseitig (Buch)
    toc=flat,                           % Flache TOC
    %acronym=totoc                       % Abkürzungsverzeichnis         im Inhaltsverzeichnis aufführen
    listof=totoc,                       % Verzeichnisse                 im Inhaltsverzeichnis aufführen
    %bibliography=totoc,                 % Literaturverzeichnis          im Inhaltsverzeichnis aufführen
    %index=totoc,                        % Index                         im Inhaltsverzeichnis aufführen
    listof=leveldown,                   % Gleitumgebungsverzeichnisse   eine Ebene tiefer
    %bibliography=leveldown,             % Literaturverzeichnis          eine Ebene tiefer
    %index=leveldown,                    % Index                         eine Ebene tiefer
    captions=tableheading,              % Beschriftung von Tabellen unterhalb ausgeben 
    final                               % Status des Dokuments (final/draft)
]{scrreprt}

% Meta-Informationen -----------------------------------------------------------
%   Definition von globalen Parametern, die im gesamten Dokument verwendet
%   werden können (z.B auf dem Deckblatt etc.).
% ------------------------------------------------------------------------------
\usepackage[latin1]{inputenc}

\newcommand{\art}{Bachelorarbeit}
\newcommand{\jahr}{2018}

\newcommand{\titel}{Titel der Bachelorarbeit}
\newcommand{\untertitel}{Untertitel der Bachelorarbeit}

\newcommand{\fachgebiet}{Fachgebiet}                        % Informatik
\newcommand{\studienbereich}{Studienbereich}                % Angewandte Informatik
\newcommand{\Vertiefungsstudium}{Schwerpunkt}               % Wirtschaftsinformatik

\newcommand{\wbh}{Name der HS/Uni}                          % Wilhelm Büchner Hochschule
\newcommand{\wbhlogo}{ba_logo/WBH_1000px}                   % .ps
\newcommand{\logo}{ba_logo/logo}                            % .ps

\newcommand{\autor}{Vorname Nachname}
\newcommand{\matrikelnr}{888888}

\newcommand{\erstgutachter}{Prof. Dr. XXX}
\newcommand{\zweitgutachter}{Dipl.-Inf. XXX}
\newcommand{\abgabedatum}{01.01.2018}

\newcommand{\ort}{Ort}
\newcommand{\datum}{01.01.2018}                          % Metadaten für die folgenden Dokumentteile
% packages ---------------------------------------------------------------------
% Die Datei packages.tex beinhaltet alle benötigten packages alsauch diese,
% welche später verfügbar sein sollten.
% ------------------------------------------------------------------------------

% Seitenstil -------------------------------------------------------------------
\usepackage[
    automark,                           % Kapitelangaben in Kopfzeile automatisch erstellen
    headsepline,                        % Trennlinie unter Kopfzeile
    ilines                              % Trennlinie linksbündig ausrichten
]{scrpage2}

\usepackage[ngerman]{babel}             % ...
\usepackage{moreverb}                   % ...
\usepackage{ae}                         % Schönere Schriften zur PDF-Erstellung

% Schrift ----------------------------------------------------------------------
\usepackage{lmodern}                    % neue deutsche Rechtschreibung zur Silbentrennung,
                                        % korrekten Benennung von Inhaltsverzeichnis und Literaturverzeichnis
\usepackage{relsize}                    % Schriftgröße relativ festlegen

% Umlaute ----------------------------------------------------------------------
\usepackage{ngerman}
\usepackage[latin1]{inputenc}
\usepackage[T1]{fontenc}                % zur korrekten Darstellung von öüä

% Zeichensätze -----------------------------------------------------------------
\usepackage{textcomp}                   % \textcopyright, \texttrademark
\usepackage[official,right]{eurosym}    % EUR-Symbol; \euro oder \EUR{XX,XX}

% Grafiken ---------------------------------------------------------------------
\usepackage[dvips,final]{graphicx}      % Einbinden von JPG-Grafiken ermöglichen
\graphicspath{{ba_logo/}}               % hier liegen die Bilder des Dokuments
\usepackage{floatflt}                   % umfließen von Bildern

% Index-Ausgabe mit \printindex ------------------------------------------------
\usepackage{makeidx}

% Zeilenabstände und Seitenränder etc. -----------------------------------------
\usepackage{setspace}                   % Definition von Zeichen-/Zeilenabständen
\usepackage{geometry}                   % Definition von Seitenrändern

% Index & Hyperlinks -----------------------------------------------------------
\usepackage{makeidx}
\usepackage[colorlinks=true, linkcolor=blue]{hyperref}
%\usepackage{hyperref}                   % genutzt für den Glossar
\usepackage{url}                        % URL verlinken

% Glossar - https://en.wikibooks.org/wiki/LaTeX/Glossary -----------------------
\usepackage{glossaries}                 % zur Erstellung eines Glossars

% Zitierweise ------------------------------------------------------------------
\usepackage[numbers]{natbib}            % Zitierweise, ohne options wird ein Fehler bei fehlendem Jahr im Literaturverzeichnis angezeigt

%\usepackage[backend=bibtex , style=numeric , bibstyle=numeric ]{biblatex}
%\usepackage[style=authoryear,backend=biber,natbib=true,sortlocale=de_DE,sorting=none]{biblatex}
%\addbibresource{literaturliste.bib} %Literatur laden
%\usepackage{csquotes}% Recommended




%\usepackage[
%    style=authoryear,           % Zitierstil
%    isbn=false,                 % ISBN nicht anzeigen, gleiches geht mit nahezu allen anderen Feldern
%    pagetracker=true,           % ebd. bei wiederholten Angaben (false=ausgeschaltet, page=Seite, spread=Doppelseite, true=automatisch)
%    maxbibnames=50,             % maximale Namen, die im Literaturverzeichnis angezeigt werden (ich wollte alle)
%    maxcitenames=3,             % maximale Namen, die im Text angezeigt werden, ab 4 wird u.a. nach den ersten Autor angezeigt
%    autocite=inline,            % regelt Aussehen für \autocite (inline=\parancite)
%    block=space,                % kleiner horizontaler Platz zwischen den Feldern
%    backref=true,               % Seiten anzeigen, auf denen die Referenz vorkommt
%    backrefstyle=three+,        % fasst Seiten zusammen, z.B. S. 2f, 6ff, 7-10
%    date=short,                 % Datumsformat
%    ]{biblatex}
%\setlength{\bibitemsep}{1em}    % Abstand zwischen den Literaturangaben
%\setlength{\bibhang}{2em}       % Einzug nach jeweils erster Zeile
%\bibliography{literature}       % Bibtex-Datei wird schon in der Preambel eingebunden

% abgelegte packages -----------------------------------------------------------
%\usepackage{typearea}           % Fehler zur Präambel-Head-Zeile umgehen


                      % Packages die benötigt werden
\include{befehle}                       % Befehle
% Seitenstil -------------------------------------------------------------------
%   Diese Datei überträgt die Einstellungen auf die folgenden Seiten
%   des Dokuments
% ------------------------------------------------------------------------------

% Satzspiegel -----------------------------------------------------------------
\setlength{\topskip}{\ht\strutbox} % behebt Warnung von geometry
\geometry{paper=a4paper,left=35mm,right=35mm,top=30mm,bottom=50mm}

% Schriftart -------------------------------------------------------------------
\usepackage{times}
%\usepackage{fontspec}
%\setmainfont{Arial}

% Zeilenabstand 1,5 Zeilen -----------------------------------------------------
\onehalfspacing

% Kopf- und Fußzeilen ----------------------------------------------------------
\pagestyle{scrheadings}
\renewcommand*{\chapterpagestyle}{scrheadings} 
% Schriftform der Kopfzeile
%\renewcommand{\headfont}{\normalfont}

\ihead{                                     % Kopfzeile
    \textbf{\large{\titel}}                 % Titel der Arbeit aus der meta.tex
    \\[5px] 
    \large{\untertitel}                     % Untertitel der Arbeit aus der meta.tex
}
\chead{} 
\ohead{\includegraphics[scale=1.0]{\logo}} 
\setlength{\headheight}{21mm}
%\setheadsepline[text]{0.4pt}                % Linie unter die Kopfzeile
%\setwidth{0pt}{textwithmarginpar}           % nur wenn außerhalb des Textes erweitert werden soll

\ifoot{\headmark}                           % Fußzeile
\cfoot{}                           
\ofoot{\pagemark}                           % Seitenzahl                    % Seitenstil für die folgenden Dokumentteile
\begin{filecontents}{bachelorarbeit_lit.bib} 
    @book{aristotle:physics, 
        author       = {Aristotle}, 
        title        = {Physics}, 
        date         = 1929, 
        translator   = {Wicksteed, P. H. and Cornford, F. M.}, 
        publisher    = {G. P. Putnam}, 
        location     = {New York}, 
        keywords     = {primary}, 
        langid       = {english}, 
        langidopts   = {variant=american}, 
        shorttitle   = {Physics}, 
        annotation   = {A \texttt{book} entry with a \texttt{translator} field}, 
    }
    @online{wikibook,
        title        = {Generating Bibliographies with biblatex and biber},
        organization = {Wikibooks},
        date         = {2016},
        urldate      = {2016-03-07},
        url          = {https://en.wikibooks.org/wiki/LaTeX/Generating_Bibliographies_with_biblatex_and_biber},
        keywords     = {untrusted},
    }
\end{filecontents}         % Erstellen aller Zitierungen *.bib
\makeindex                              % Index

\begin{document}
    \pagenumbering{Roman}                               % I (Deckblatt), II (Eidesstattliche Erklärung), III (Danksagung), IV (Abbildungsverzeichnis), ...
    % Deckblatt der Bachelorarbeit -------------------------------------------------
% 
% ------------------------------------------------------------------------------
\thispagestyle{plain}
\begin{titlepage}
    \begin{center}
        \huge{\textbf{\titel}}\\[1.5ex]
        \LARGE{\textbf{\untertitel}}\\[6ex]
        \LARGE{\textbf{\art}}\\[1.5ex]
        \Large{im Fachgebiet \fachgebiet}\\[12ex]

        \includegraphics[scale=0.2]{\wbhlogo}\\[6ex]

        \normalsize
        \begin{tabular}{p{5.4cm}p{6cm}l}\\
            vorgelegt von:  & \quad \autor\\[1.2ex]
            Studienbereich: & \quad \studienbereich\\[1.2ex]
            Matrikelnummer: & \quad \matrikelnr\\[1.2ex]
            Erstgutachter:  & \quad \erstgutachter\\[1.2ex]
            Zweitgutachter: & \quad \zweitgutachter\\[3ex]
        \end{tabular}

        \copyright\ \jahr\\[6ex]

    \end{center}
    \singlespacing
    \small
    \noindent Dieses Werk einschlie"slich seiner Teile ist \textbf{urheberrechtlich gesch"utzt}. Jede Verwertung au"serhalb der engen Grenzen des Urheberrechtgesetzes ist ohne Zustimmung des Autors unzul"assig und strafbar. Das gilt insbesondere f"ur Vervielf"altigungen, "Ubersetzungen, Mikroverfilmungen sowie die Einspeicherung und Verarbeitung in elektronischen Systemen.
\end{titlepage}                 % Deckblatt der Bachelorarbeit aufnehmen
    \section*{Danksagung}
\label{sec:Danksagung}
An dieser Stelle m"ochte ich mich bei all denjenigen bedanken, die mich w"ahrend der Anfertigung dieser Masterarbeit unterst"utzt und motiviert haben.
Zuerst geb"uhrt mein Dank Frau Prof. XY, die meine Masterarbeit betreut und begutachtet hat. F"ur die hilfreichen Anregungen und die konstruktive Kritik bei der Erstellung dieser Arbeit m"ochte ich mich herzlich bedanken.
Ebenfalls m"ochte ich mich bei meinen Kommilitonen XY und XY bedanken, die mir mit viel Geduld, Interesse und Hilfsbereitschaft zur Seite standen. Bedanken m"ochte ich mich f"ur die zahlreichen interessanten Debatten und Ideen, die ma"sgeblich dazu beigetragen haben, dass diese Masterarbeit in dieser Form vorliegt.
Ein besonderer Dank gilt allen Teilnehmern und Teilnehmerinnen meiner Befragung, ohne die diese Arbeit nicht h"atte entstehen k"onnen. Mein Dank gilt ihrer Informationsbereitschaft und ihren interessanten Beitr"agen und Antworten auf meine Fragen.
Meinen Freunden XY, XY und XY und meinem Freund XY danke ich besonders f"ur den starken emotionalen R"uckhalt "uber die Dauer meines gesamten Studiums.
Abschlie"send m"ochte ich mich bei meinen Eltern XY und XY bedanken, die mir mein Studium durch ihre Unterst"utzung erm"oglicht haben und stets ein offenes Ohr f"ur meine Sorgen hatten.
\\[15px]
\ort, den \datum
\\[30px]
\begin{tabular}{@{}l@{}}\hline
    \autor
\end{tabular}

\newpage               % Danksagung der Bachelorarbeit aufnehmen
    \markboth{Inhaltsverzeichnis}{Inhaltsverzeichnis}
% \section*{Inhaltsverzeichnis}
\label{sec:Inhaltsverzeichnis}
\tableofcontents                                    % Inhaltsverzeichnis erstellen

\newpage                     % Index
    \section{Abk\"urzungsverzeichnis}
\begin{acronym}
    \acro{Kuerzel}[KF]{Langform}
    \acro{EU}[EU]{Europ\"aische Union}
\end{acronym}                 % Abkürzungsverzeichnis
    \markboth{Abbildungsverzeichnis}{Abbildungsverzeichnis}
% \section*{Abbildungsverzeichnis}
\label{sec:Abbildungsverzeichnis}
\listoffigures                              % Abbildungsverzeichnis erstellen
                  % Abbildungsverzeichnis
    \markboth{Tabellenverzeichnis}{Tabellenverzeichnis}
% \section*{Tabellenverzeichnis} 
\listoftables
                    % Tabellenverzeichnis
    \pagenumbering{arabic}                              % 1, 2, 3, 4,...
    \markboth{Erstes Kapitel}{Erstes Kapitel}
\chapter{Kapitel mit Testelementen - LaTeX}                                         % 1     Kapitel mit Testelementen - LaTeX
\section{Zitierungen}                                                               % 1.1   Zitierungen
Das Zitieren im Harvard-Stil ist vor allem bei Sozial- und
Naturwissenschaftlern beliebt. Ein Hauptunterschied zu anderen Zitierweisen
ist dabei der Verzicht auf Fu\ss noten. Hei\ss t, dass du mitunter mehr Platz
pro Seite für deinen Flie\ss text hast - ob das nun ein Fluch oder Segen ist,
\"uberlassen wir deiner Vorliebe. Der Verzicht auf Fu\ss noten im
Harvard-Zitieren zieht nat\"urlich eines nach sich: Den Literaturhinweis
musst du in deinen Text einflie\ss en lassen. Doch auch da zieht die
amerikanische Elite-Universit\"at die Sparsamkeit der Vollst\"andigkeit vor.
An geeigneter Stelle verweist du n\"amlich nur auf den Nachnamen der Autorin
oder des Autors, das Erscheinungsjahr und die Seite, von der du zitierst.\par

EIn Zitat kann dabei durch eine hochgestellte Fu\ss note
\footfullcite{aristotle:physics} oder durch einen Literaturhinweis im Sinne der
Harvard-Zitierweise.                                                                % footfullcite (ausführlich) oder \footcite (kurz)

\subsection{Direktes Zitat}                                                         % 1.1.1 Direktes Zitat
Die Funktion der Ehe ist als
"Grundstein der Familie und als Zentrum der geistig-seelischen Entwicklung des
Menschen und andererseits als \"okonomisch-politischer Nutzfaktor"
\autocite{aristotle:physics} zu betrachten.                                         % Variante (ohne Fu\ss zeile, Name, Jahr)

\subsection{Indirektes Zitat}
Die Architextualit\"at besch\"aftigt sich damit, wie die Zuordnung eines Textes
zu einer Kategorie stattfindet und einer Gattung zugeordnet wird
\autocite[9]{aristotle:physics}.                                                    % Variante (ohne Fu\ss zeile, Name, Jahr, Seite)

\subsection{Auslassungen}
Die Funktion der Ehe ist als "Grundstein der Familie und als Zentrum der
geistig-seelischen Entwicklung des Menschen [...]"
\autocite[9]{aristotle:physics} zu betrachten.                                      % Variante (ohne Fu\ss zeile, Name, Jahr, Seite)

\subsection{Ver\"anderungen}
Die Funktionen von Ehen sind die "Grundstein[e] der Familie und als Zentrum
der geistig-seelischen Entwicklung des Menschen und andererseits als
\"okonomisch-politischer Nutzfaktor" \autocite[9]{aristotle:physics} zu betrachten. % Variante (ohne Fu\ss zeile, Name, Jahr, Seite)

\subsection{Falschschreibungen}
Befinden sich Tipp-, Grammatik- oder \"ahnliche Fehler im Original,
\"ubernimmst du diese zwar, kennzeichnest dies jedoch durch ein [sic!]
hinter dem jeweils fehlerhaften Wort. Aus Gr\"unden der Lesbarkeit
wird auf ein [sic!] verzichtet, wenn Textstellen zitiert werden,
die noch in alter Rechtschreibung geschrieben sind.                                 % Fehler im Orginal

\section{Indexeintr\"age}
Affe\index{Affe}
Bieber\index{Bieber}                                                                % Indexverweis durch \index{Bez}


\section{Abk\"urzungen}                                                             
Das Begrifflichkeit \ac{EU} wird in erster Verwendung auch in der Langform
dargestellt. Wird das gleiche K\"urzel (\ac{EU}) erneut aufgerufen,
so erfolgt lediglich die Kurzform.                                                  % Abk\"urzungen durch \ac{Bez}

\newpage

\section{Abbildungen}
\begin{figure}[h!]                                                                  % Option [h!] platziert das Bild fix
    \centering                                                                      % zentrieren der Abbildung
    \includegraphics[width=.9\linewidth]{ba_abbildungen/testimage}                  % .ps, width=10cm, scale=0.5, width=.5\linewidth
    \captionof{figure}{Eine Beispielabbildung}                                      % \Beschriftung der Abbildung
\end{figure}

\section{Tabellen}
\begin{table}[h!]                                                                   % Option [h!] platziert die Tabelle fix
    \centering                                                                      % zentrieren der Tabelle
    \begin{tabular}{l|c|r}\hline\hline                                              % Definition der Tabelle (spalten, links-, zentriert- oder rechts-bündig)
        links & zentriert & rechts \\\hline A & B & B\\\cline{2-2}                  % Verbinden von Zeilen
        1 & 2 & 3\\\hline
    \end{tabular}
    \captionof{table}{Tabellenbeschreibung}                                         % Beschriftung der Tabelle
\end{table} 

\newpage                       % Teilung in Einleitung, Hauptteil und Schluss/Fazit
    \appendix
\section*{Anhang}
\label{sec:Anhang}
Der Anhang
    \markboth{Literaturverzeichnis}{Literaturverzeichnis}
\printbibliography

\newpage   % Literaturverzeichnis
    \markboth{Index}{Index}
% \section*{Index}
\label{sec:Index}
\printindex                                         % Index erstellen

\newpage                                  % Index
    \thispagestyle{empty} 
\section*{Eidesstattliche Erkl"arung}
%\label{Eidesstattliche Erkl"arung}
Hiermit erkl"are ich, dass ich die vorliegende Arbeit eigenst"andig und ohne fremde Hilfe angefertigt habe. Textpassagen, die w"ortlich oder dem Sinn nach auf Publikationen oder Vortr"agen anderer Autoren beruhen, sind als solche kenntlich gemacht.Die Arbeit wurde bisher keiner anderen Pr"ufungsbeh"orde vorgelegt und auch noch nicht ver"offentlicht.
\\[15px]
\ort, den \datum
\\[30px]
\begin{tabular}{@{}l@{}}\hline
    \autor
\end{tabular}     % Eidesstattliche Erklärung der Bachelorarbeit aufnehmen
\end{document}